\chapter{Introduction}% Main chapter title
\label{Chapter1} % For referencing the chapter elsewhere, use \ref{Chapter1}

\section{Welcome}

Welcome to the manual of the Uni Leipzig Bioinf Thesis class. This class gives
a solid backbone for your thesis writing at the \quo{Institut für Informatik}
at the \quo{Universität Leipzig}. It can be used as it is without any further
configuration. However, many things can be configured. So hopefully it fits
your needs and your idea of design.

This template is designed to fulfill the requirements of the
\quo{Promotionsordnung Fakultät für Mathematik und Informatik} as it is in 2018
\citep{web:promotion,web:promotion2}. So in case of changes in the regulation
some parts of this work may not longer be suitable.

This class has many class options and also introduces some new commands that
can be used. So when you really want to use it you should at least read
\autoref{AppendixA} and \autoref{AppendixB}. However, I would recommend you to
read the complete manual so you get an overview how things work and how things
should be done. Also, the short overviews in the Appendix reflect not every
single aspect of the commands and options.

This manual is not an introduction to \LaTeX. The authors assume that you
already typeset a scientific work with \LaTeX{} before and that you do not
need an introduction to the basic commands. The typesetting of mathematical
formulas is not covered, either.

The use of this class is free and you can change things by yourself and
redistribute it if you like. However, if you use it and it makes your life
easier, I would appreciate it when you mention me in your acknowledgments.
Something short like: \quo{Thanks to Fabian Gärtner for the latex class of
this thesis.}

If something does not work as expected, something needs an update because of
some changes or you miss an option to modify the template feel free to write
me an email: \email{fabian@bioinf.uni-leipzig.de}.  I will see if I can help
you.

\subsection{Disclaimer}
Before we start, here is a little warning. At the time I made this template, I
was a PhD student.  I have tried to be as accurate as possible in everything
that is written in this manual,  and also in choosing its default options.
However, some people might disagree with me and there is not just one correct
answer but rather several possible options. So trust your supervisor more than
anything I have done here. Again if you get some feedback that something has
to change I would appreciate hearing from you such that I can include changes
in order to keep the template updated.


\subsection{This Manual}
This manual was itself created with the given
\pkg{UniLeipzigBioinfThesis} class. In the original download, the source
files contain the code of this manual so that you have direct examples how
things are done in this class. Of course, you will delete it in order to write
your own content in the files. However, before you do it you may have a look
at the content and keep the pdf file in a separate location such that you can
reread the manual. You can also download it again in order to see the original
sources.

The only difference between this source and the source of \file{manual.pdf}
is the class option \clsopt{minted} (see \autoref{sub:codehigh}) and
\clsopt{zusammenpath} (see \autoref{sec:germanabstract}) are set different.
This is done to make it compile and a offer good starting point without any
further changes.

\section{Double-sided Print}
This style only supports double-sided layout for the thesis. This is not a
lack of an option, this is done purposely.  When reading on a screen, it makes
more or less no difference if the document is single-sided or double-sided.
The difference is in the way it is printed.  A double-sided print is done such
that every sheet of paper is printed on both sides. This has only advantages.
It is closer to normal books and therefore more natural to read. It makes the
thesis thinner so that it is easier to handle. It saves you money because you
print on and tie less pages. The last and best reason is that it reduces the
waste of paper and makes the world a better place. I believe this is the
reason why we all write a thesis, to make the world a better place.


\section{File Structure}
This template splits the content of the thesis into various different files.
This helps you to keep your work organized and allows you to only compile the
parts that you currently edit. It also enables the automatic generation of
stand-alone documents from several parts of the work (e.\,g.\ the abstract),
which are formally required for the Dekanat when you hand in your thesis.
Other files provide the functionality of the class itself and need not be
changed.  Another set of files helps you to compile the thesis and store your
personal information.

The thesis itself is compiled from the file \file{main.tex} in the root
directory. In this file, the thesis class is loaded and basic configuration is
done. Especially, every path to another file can be changed here so that all
files can be moved or renamed if necessary. More details are given in
\autoref{sec:main}. Then, all the content from the subfolders is included and
compiled into a single document. The class itself is taking care of creating
the front matter and back matter for you. In addition to the main file, one
tex file for each stand-alone document created from the thesis content is
provided:
\begin{description}
    \item[\file{abstract\_sa.tex}:]
        The stand-alone abstract.
    \item[\file{curriculum\_scientiae\_sa.tex}:]
        The curriculum outlining your scientific career.
    \item[\file{publications\_sa.tex}:]
        The list of all your scientific publications.
    \item[\file{selbststaendigkeitserklaerung\_sa.tex}:]
        The declaration that you are the original author of all your thesis'
        content and that you have properly marked all citations.
\end{description}

The classes for the thesis and the stand-alone documents are defined in the
files \file{UniLeipzigBioinfThesis.cls} and
\file{UniLeipzigBioinfThesisSA.cls}, respectively. A style file,
\file{UniLeipzigBioinfThesisCommon.sty}, contains code that is shared by
both classes.  These classes make all commands and styles available to the tex
files. You should not edit these files unless you want to add a new feature or
have to fix a bug.

The files \file{Makefile} and \file{latexmkrc} help you to compile the
documents, as explained in \autoref{sec:compilation}.

The folders have thematically different content. The \file{frontmatter}
folder contains files that are used in the front matter (abstract,
acknowledgments,\dots), the \file{chapters} folder contains a single tex
file for each chapter of the main text. The \file{graphics} folder contains
the graphics for the figures of the thesis.  The \file{appendices} folder
contains a single tex file for each appendix.  The last folder is
\file{backmatter} which contains files that are used in the back matter
(\eg the Curriculum Scientiae).  The \file{misc} folder contains two data
files and a tex file. \file{literature.bib} is the \prog{BibTeX} file that
contains all works that you cited from. Your personal data and thesis
information like title, keywords etc., is stored in \file{thesis\_info.tex}.
File \file{custom\_defs.tex} is for all commands you may want to define
yourself, as well as for loading other packages you may want to use. It comes
with some commands already that you can use or change if you like.

Finally, \file{manual.pdf} contains a copy of this manual. You may want to
keep it for future reference.


\section{How to Compile It}%
\label{sec:compilation}

The short version: run \cmd{make} to compile everything, or \cmd{make
main} to only compile the thesis. Run \cmd{make watch} to continuously
compile the thesis as you write. If you have problems or want more
information, please read on.

As of release v2.0 of this class, the build system has been completely
rewritten. It now relies on \prog{latexmk}, a dedicated tool for compiling
\LaTeX{} documents that is taking care of all required steps, including the
execution of \prog{bibtex}, \prog{biber}, \prog{makeindex},
\prog{pdflatex} etc., in the right order and as often as necessary. It is
also more efficient than manually calling all these tools, because it
determines which files have changed and regenerates the output as necessary.
Furthermore, it offers file monitoring to automatically recompile the
document as soon as any changes to the source code are saved, as well as
cleaning procedures. Still, a \file{Makefile} is provided to allow calling
\prog{latexmk} with a single command without having to remember all the
options and paths. For more information on how to compile using the provided
\file{Makefile}, refer to \autoref{sub:make}.

\subsection{Configuring \prog{latexmk}}
With this template comes the file \file{latexmkrc}, which already contains
the required configuration for \prog{latexmk}. Separating the configuration
from the \file{Makefile} has the advantage that \prog{latexmk} can also be
called via any third-party \LaTeX{} development environment or editor without
further configuration. For example, the compilation feature of the
\prog{vimtex} plugin for the editor \prog{vim} works out of the box with
the provided files.


\subsection{Makefile}%
\label{sub:make}

The provided \file{Makefile} is divided into the sections \quo{Config} and
\quo{Rules}. Usually, you don't have to touch the rules at all, and
changing the configuration variables is only necessary when you rename the
main or stand-alone tex files.

To compile the thesis and all other stand-alone document, run \cmd{make
all} or simply \cmd{make} in the thesis directory. \cmd{all} is called
the (default) \emph{target}, \ie the object that is to be made. The
\file{Makefile} also defines several other useful targets:
%
\begin{description}
    \item[main:] Compiles the main file, \ie the thesis. The main file is
        by definition the first file listed in the variable
        \cmd{TEX\_FILES}. Running \cmd{make main} is the quickest
        way to compile only the thesis, omitting the stand-alone documents.
    \item[\emph{filename}.pdf:] Compiles (only) \file{filename.pdf} from
        source \file{filename.tex} for any file name.
    \item[clean:] Removes \emph{all} generated files, including the pdf files.
    \item[watch:] Watches for changes to any of the sources of the main file
        and re-compiles them as soon as any are detected.
    \item[watch-filename.tex:] Like \cmd{watch}, but watch
        \file{filename.tex} instead of the main file.
\end{description}
%
You can also set the path to \prog{latexmk} and additional command line
arguments in the \file{Makefile}. A few additional targets are provided for
development purposes:
%
\begin{description}
    \item[dist:] Make a tarball (\file{.tar.gz} archive) of the thesis
        directory. Target \cmd{clean} is executed first. The file name can
        be set by changing the value of variable \cmd{DIST\_NAME}. Set
        \cmd{DIST\_EXCLUDE} to prevent certain files or directories from
        being included in the archive.
    \item[changes:] Update \file{CHANGES.md} by converting
        \file{appendices/changes.tex} (\ie the version history) to MarkDown
        using \prog{Pandoc}.
    \item[html:] Generate HTML versions of \file{README.md} and
        \file{appendices/changes.tex} to update the website of the thesis
        template.\footnote{\url{%
        https://www.bioinf.uni-leipzig.de/intern/phd-thesis-class/}}
\end{description}

\subsection{Dependencies and Compiling at the Bioinf Institute}
To compile this template, you obviously need a \LaTeX{} distribution. Under
Linux, this is usually \prog{TeXLive}, which is available in the official
package repositories of any major Linux distribution. However, as the template
has many features, it also depends on many packages. The easiest way to
provide them all is a full installation of all \prog{TeXLive} packages. In
Fedora 36, this is done by installing package \pkg{texlive-scheme-full}
using \cmd{dnf}. In Ubuntu 22.04, install \pkg{texlive-full} and
\pkg{python3-pygments} with \cmd{apt}.  Also make sure that
\cmd{make} is available if you want to use the \file{Makefile}.  As of
2022, the computers at the Bioinformatics institute run varying versions of
Fedora. If the required packages are not installed on your machine yet, please
send a short email to the system administrator --- Jens will gladly help you
out.

Though this template has not been tested on Windows, you may have luck with
\prog{MikTeX} as well.


\section{The Main File}%
\label{sec:main}

The main file can be divided into two parts: the preamble and the document
body. In the preamble the class is loaded with the desired options. In the
document body, the actual content of the thesis is included from other tex
files.

\subsection{Preamble}
The first thing that is done in the preamble is loading the required classes.
This is done with \verb|\documentclass[]{UniLeipzigBioinfThesis}|.  Class
options can be added in the square brackets.  The possible options are spread
in the manual but all are marked with a margin note.

After this, you need to attach one or more bib files to you thesis this is
explained in \autoref{sec:bib}. You can also defined additional commands and
load other packages here if necessary.

\subsection{Document Body}%
\label{sub:body}

The document starts with the \verb!\makefrontmatter! command. This command
creates the complete front matter of the thesis. This includes the title page,
the dedication, the bibliographic description, the abstract, the German
abstract, the acknowledgments and \abb{ToC}. Make sure that your personal data
in file \file{thesis\_info.tex} is correct as it is used to generate the
front matter.

The main matter of the thesis start with the call of the
\verb|\beginmainmatter| command. The command starts the page numbering with
Arabic numerals. The main matter is composed of chapters. Every chapter is in
a separated file. This files then are included in the main file with the
\verb|\include| command. For more details on chapters, have a look at
\autoref{sec:chapter}. The only other elements in the main matter are the
parts.  Parts are described in more detail in \autoref{cmd:part}.

The next part is the appendix. To start this part use the
\verb|\beginappendix| command. It changes the chapter behavior and adds an
entry to the \abb{ToC}.  Every appendix is in an extra file. To add these
files, use the \verb|\include| command again. For more information, see
\autoref{sec:app}. The last command is then \verb|\makebackmatter|. This
command creates the back matter that includes the List of Xs, bibliography,
curriculum scientiae and the declaration.

\section{Base Settings}
\subsection{Margins}%
\label{opt:margins}

By default, the template uses relatively wide margins as defined by the
default settings of the \emph{Memoir} class. Though this may seem unusual
first, especially when you are used to a typical Word document, but it is
generally considered the more aesthetic choice by most typographers. It also
reduces the eye movement  compared to longer lines and thus can be read more
easily. However, when you have much to say, this can make your thesis very
large and heavy when printing. Thus, it could make sense to reduce the margin
size. This of course reduces the space for margin notes and I would not use
them with small margins. You can change to smaller margins with the class
option \clsopt{smallmargins}\indexdef{smallmargins}.

\subsection{Font}%
\label{opt:font}

The default font that is used in the tex files is Latin modern with serifs.
This is different to usual latex documents that use Computer modern with
serifs as default font, but the result is usually much better.   You can
change the font with the class option \clsopt{font}\indexdef{font}.  The
possibilities are shown in Table~\ref{tab:fontoption}.

\begin{SCtable}
  \captiont{Font values.}{The values that the class option \quo{font} can have.}%
  \label{tab:fontoption}
    \begin{tabular}{l l}
        \toprule
        \tblhead{Value} & \tblhead{Font used}\\
        \midrule
        lm  & Latin modern without serifs \\
        lms & Latin modern with serifs (default)\\
        cm  & Computer modern without serifs\\
        cms & Computer modern with serifs\\
        \bottomrule
    \end{tabular}
\end{SCtable}


\subsection{Chapter Design}%
\label{sub:chapterdesign}

There are many different possibilities for chapter page design. To change this,
use the class option \clsopt{chapterstyle}.\indexdef{chapterstyle}.  An example
of the possible designs is shown in
Figures~\ref{fig:veelo}--\ref{fig:wilsondob}. However, you could use any design
of the \emph{Memoir} class shown in \citep[P:383-396]{wilson2018memoir}. The
default value is the \emph{veelo} style, which has a mini \abb{ToC} within
each start page of a new chapter.


\begin{figure}
    \includegraphics[scale=0.6]{veelo.pdf}
    \captiont{Veelo page style example.}{%
        An example how the veelo page style look with a mini \abb{ToC}}%
    \label{fig:veelo}
\end{figure}

\begin{figure}
    \includegraphics[scale=0.6]{default.pdf}
    \captiont{Default page style example.}{%
        An example how the default page style look with a mini \abb{ToC}}

\end{figure}

\begin{figure}
    \includegraphics[scale=0.6]{bianchi.pdf}
    \captiont{Bianchi page style example.}{%
        An example how the bianchi page style look with a mini \abb{ToC}}
\end{figure}

\begin{figure}
    \includegraphics[scale=0.6]{dash.pdf}
    \captiont{Dash page style example.}{%
        An example how the dash page style look with a mini \abb{ToC}}
\end{figure}

\begin{figure}
    \includegraphics[scale=0.6]{ell.pdf}
    \captiont{Ell page style example.}{%
        An example how the ell page style look with a mini \abb{ToC}}
\end{figure}

\begin{figure}
    \includegraphics[scale=0.6]{wilsondob.pdf}
    \captiont{Wilsondob page style example.}{%
        An example how the wilsondob page style look with a mini \abb{ToC}}%
    \label{fig:wilsondob}
\end{figure}


\section{Common Typographic Mistakes}

While \LaTeX{} is designed to generate documents with a high typographic
quality, there are a few common pitfalls that authors have to be aware of to
avoid unexpected or unpleasant results. Here, some of the most relevant tricky
details to keep in mind are described.

\subsection{Logical Markup}
One of the most important rules to achieve a consistent typography in your \LaTeX{}
documents is to consequently use \emph{logical markup} like \verb!\emph! (for
emphasis) or \verb!\email! (for an email address) instead of directly styling
the text with low-level directives like, \eg, \verb!\textit! (for italic text)
or \verb!\tiny! (for small a font size). Logical markup means that you apply a
command to a piece of text that describes \emph{what it is}, and not what it
looks like. For example, a file name is often typeset in typewriter font.
However, instead of directly using \verb!\texttt{the_name}!  to achieve this,
it is better to define a command \verb!\file{}! for this purpose and use
\verb!\file{the_name}!. This has multiple benefits: (1) All file names using
this command will look identically, and you do not need to remember the exact
style. (2) You can easily change the style later by simply changing definition
of the logical markup command. (3) A custom command may be shorter than the
basic \LaTeX{} commands, especially if you have to use multiple ones to
achieve the desired look. (4) You can simply find all file names etc.\ in the
document by searching for the respective logical markup command. In case you
want to get rid of a logical markup command later, you can usually simply
replace all occurrences using a regular expression substitution.

To define your logical markup commands, use the \verb!\newcommand! command.
User-defined commands should be defined in file \file{misc/custom\_defs.tex}.
It already contains pre-defined commands for the most common text elements in
our fields, \eg \verb!\species! for species names, \verb!\prog! for computer
program names, and \verb!\cmd! for shell commands and code. For a complete
list, refer to \autoref{tab:logical_markup}. Of course, you can extend the
list or change the styles if you like.

\begin{SCtable}
    \centering
    \caption{%
        List of the default logical markup commands, which are defined in
        \file{misc/custom\_defs.tex}.
    }%
    \label{tab:logical_markup}
    \begin{tabular}{l l}
        \toprule
        \tblhead{Command} & \tblhead{Description} \\
        \midrule
        \verb!\species! & species name \\
        \verb!\seq!     & DNA/RNA sequence \\
        \verb!\file!    & file name \\
        \verb!\cmd!     & code or command \\
        \verb!\prog!    & program name \\
        \verb!\email!   & email address \\
        \verb!\quo!     & quotation marks \\
        \verb!\tblhead! & table headers \\
        \verb!\pkg!     & software package/module \\
        \bottomrule
    \end{tabular}
\end{SCtable}

\subsection{Spaces}
Inter-word spaces are a common source of typographic mistakes.  This is mainly
because, in \LaTeX, not every space has the same size. Most notably, it
increases the space after a period (\quo{\,.\,}), which is usually good as it
helps to visually separate different sentences. However, the period is also
used to mark abbreviations, which leads to wrong spacing as, \eg, in Mr. A.
Brown. The solution is to escape the following space, as in Mr.\ A.\ Brown
(using the code \verb!Mr.\ A.\ Brown!).

When defining a command which ends in a period, you do not know whether there
will be a space after it or not. In such cases, the sequence \verb!\@! can be
used to disable the increased size of a space that will possibly follow. For
example, \verb!\newcommand{\dr}{Dr.\@}! defines the command \verb!\dr!, which
is then displayed as Dr.\@ with correct spacing.

Beside regular spaces, \LaTeX{} also offers the command \verb!\,! to generate
half spaces. They are useful to group digits like 1\,500\,000 or to typeset
abbreviations consisting of multiple words like w.\,r.\,t.\ (\quo{with respect
to}), which looks much better than both w.\ r.\ t.\ (using full spaces) or
w.r.t.\ (using no spaces).

Another variety of the normal space is the \emph{unbreakable} space, written
as tilde character (\verb!~!). It has the same size as a regular space, but
prevents the line from being broken at this specific point. This is useful for
groups of text elements that belong together and should not be separated, for
example a title and a name like Mr.~Brown, or to prevent that abbreviations
followed by a period end a line and thus confuse the reader as it looks as
though the period ends the sentence.

A final recommendation is related to the somewhat surprising behaviour of
\LaTeX{} when it comes to spaces following commands. Using a command without
any arguments like \verb!\dr! defined above \quo{swallows} the following
space, \ie it will print Dr.and no space after it like so. To fix this, either
use \verb!\dr{}!, or use the \verb!\xspace! command from the \pkg{xspace}
package in the definition: \verb!\newcommand{\dr}{Dr.\@\xspace}!. It will
expand to a space only when needed and not, \eg, when used within parentheses.
The pre-defined markup commands in \file{misc/custom\_defs.tex} showcase the
use of \verb!\xspace!.

\subsection{Quotation Marks}

In short: always use the \verb!\quo{}! command to quote your text like
\quo{this}.

While using proper quotation marks like \quo{these} seems to be easy at a
first glance, it turns out to be a complex matter. First, one should realize
that the double quote character " looks horrid and should never be used in a
\LaTeX{} text. Instead, two backticks (\verb!``!) and two apostrophes
\verb!''! are used to create correct double quotation marks as in the first
sentence. Note how the opening and closing marks differ slightly. To make
things worse, quotation marks are language-specific, so in the German
Zusammenfassung, you would for instance use \glqq{}diese Art von\grqq{}
quotation marks.  Since this is easy to get wrong, this thesis class properly
sets up the language settings for the respective parts and provides the
command \verb!\enquote{text}! from the package \pkg{csquotes}, as well as the
shorthand \verb!\quo{text}!. Using this command is not only convenient and
guaranteed to produce correct results, it also handles edge cases like inner
and outer quotes: \quo{\quo{Hello,} he said.}

\subsection{Dashes}

Texts contain a wide variety of dashes---\ie horizontal lines---that have
different lengths depending on their function. To insert partial sentences
that could also be surrounded by commas or parentheses, or to append them as
if following a semicolon, the \emph{em-dash} (German Geviertstrich) is used.
It is named after the widest letter M.  In \LaTeX{}, it is typeset using three
dash characters (\verb!---!).  The first sentence of this section demonstrates
its usage. The shorter \emph{en-dash} (German Halbgeviertstrich), has half the
size of an em-dash. It can be typeset using two dashes (\verb!--!). It is used
to indicate ranges (\eg 10--15 minutes), or for compound nouns as in
Watson--Crick base pair or protein--protein interaction. The shortest dash is
the hyphen (German Viertelgeviertstrich), written simply as \verb!-! in the
code. It is used to hyphenate words at the end of lines, or for compound
adjectives like \quo{fast-growing} and sometimes for prefixes like in
\quo{re-evaluate}.

To get correctly sized minus signs, always set negative numbers in a math
environment: $-5$ (code: \verb!$-5$!) and not -5.

\subsection{Position Specifiers of Floats}

Do not overuse position specifiers of floats environments. For example, the
figure environment allows to pass the specifiers h (\emph{here}), t
(\emph{top}), and b (\emph{bottom}) as in \verb!\begin{figure}[htb]!. Do not
use them unless you have a good reason. Do not add them to all of your figures
by default. \LaTeX{} usually does a good job finding a suitable position on
its own. If you feel you have to add specifiers, do it when you are close to
finishing that chapter, when you have already included all other figures, or
you may get surprising results.

\subsection{Tables}

Instead of using \verb!\hrule! to draw horizontal lines in tables, use
\verb!\toprule!, \verb!\midrule!, and \verb!\bottomrule! as in the tables of
this manual, \eg~\autoref{tab:logical_markup}. Do not overuse these commands
as too many lines will only clutter your table, but not benefit its
readability.


% End of chapters/chapter_1.tex
